\documentclass[11pt]{article}
\usepackage[margin=1in]{geometry}
\usepackage{amsmath, amssymb}
\usepackage{siunitx}
\usepackage{graphicx}
\usepackage{booktabs}
\usepackage{hyperref}
\hypersetup{colorlinks=true, linkcolor=black, citecolor=black, urlcolor=blue}

\title{Simulating Mechanical Energy Transfer in Rowing\large\\A Simplified Fluid‑Dynamics and Mechanics Model}
\author{Gabriel Cirker}
\date{\today}

\begin{document}
\maketitle

\begin{abstract}
Rowing shells are propelled by intermittent strokes of multiple rowers.  During
the drive phase each rower applies a force through the oar blade, generating
thrust on the boat, whereas during the recovery the boat coasts and decelerates
under hydrodynamic resistance.  Understanding the conversion of rower effort
into boat motion requires a coupled description of the stroke cycle and the
drag force experienced by the hull.  This study formulates a simple one‑dimensional
model of a racing shell that balances a prescribed thrust profile against a
quadratic drag law.  Newton’s second law is integrated numerically over
multiple strokes to evaluate boat speed, energy input, drag losses and overall
efficiency.  The model reproduces realistic speeds and demonstrates that
roughly 6~\% of the mechanical work delivered by the rowers appears as an
increase in kinetic energy, the remainder being dissipated by drag.  The
framework illustrates how modest parameter changes influence performance and
provides an accessible platform for exploring rowing dynamics.
\end{abstract}

\section{Introduction}

Rowing is both a team sport and a mechanical system: multiple athletes work
together to convert metabolic energy into forward motion of a lightweight hull.
The oar blades interact with water and air, imparting momentum to the fluid
and in turn providing thrust to the boat.  Meanwhile the hull experiences
hydrodynamic drag that opposes motion.  For an eight–person shell with a
coxswain (an “8+”), typical component masses are about \SI{655}{\kilogram}
for the rowers, \SI{89}{\kilogram} for the shell and \SI{22.5}{\kilogram}
for the oars【779142687997843†L270-L287】, giving a total mass near
\SI{766}{\kilogram}.  The drag force on such a shell scales approximately
with the square of its speed; experimental studies have expressed the drag as
\(F\_d = C v^n\) with exponent between 1.5 and 2【779142687997843†L256-L261】.
In practice, the constant \(C\) encapsulates both the hydrodynamic and
aerodynamic contributions and is determined empirically; values around
9--14 \si{N\,s^2/m^2} have been measured for racing shells【779142687997843†L388-L450】.
The fundamental drag equation from fluid dynamics is
\(F\_d = \tfrac{1}{2}\,\rho u^2 c\_d A\), where \(\rho\) is the fluid
density, \(u\) is the relative flow speed, \(c\_d\) is a dimensionless drag
coefficient and \(A\) is a reference area【443245353880283†L120-L134】.  When
combined into a single constant, this expression yields the simple quadratic
law used here.

This paper develops a minimalist time‑domain model of mechanical energy
transfer in rowing.  By prescribing a periodic thrust and balancing it
against quadratic drag, the model predicts speed oscillations over the stroke
cycle, quantifies energy losses and estimates the efficiency of the stroke.

\section{Mathematical Model}

\subsection{Equations of motion}

Consider a boat of mass \(m\) moving in one dimension with velocity
\(v(t)\).  The net force on the boat is the difference between the thrust
applied by the rowers, \(F\_{\text{thrust}}(t)\), and the hydrodynamic
drag, \(F\_d\).  Newton’s second law gives
\begin{equation}
  m\,\frac{\mathrm{d}v}{\mathrm{d}t} = F\_{\text{thrust}}(t) - F\_d(v).
  \label{eq:newton}
\end{equation}
The drag force is approximated by the quadratic law
\begin{equation}
  F\_d(v) = k\_\mathrm{d} \, v^2,
  \label{eq:drag}
\end{equation}
where \(k\_\mathrm{d}\) is an empirical constant summarising the factors
\(\tfrac{1}{2}\rho c\_d A\) in the standard drag equation【443245353880283†L120-L134】.  The exponent of 2 reflects the quadratic
dependence of drag on speed【443245353880283†L182-L185】.  Experimental
measurements of 8+ shells suggest \(k\_\mathrm{d}\) values of order
\numrange{9}{14} \si{N\,s^2/m^2}【779142687997843†L388-L450】; to achieve
realistic steady‑state speeds in the present model we adopt
\(k\_\mathrm{d}=\SI{40}{N\,s^2/m^2}\).

\subsection{Thrust model}

Rowers deliver power in pulses.  During the drive phase of each stroke the
oar blades engage the water and the crew pushes hard against the footplates,
producing thrust.  During the recovery the blades are feathered and the boat
coasts.  To capture this behaviour, the thrust is modelled as a square wave

\begin{equation}
  F\_{\text{thrust}}(t) =
    \begin{cases}
      F\_\text{drive}, & 0 \le (t \bmod T) < \tau,\\
      0, & \tau \le (t \bmod T) < T,
    \end{cases}
  \label{eq:thrust}
\end{equation}

where \(T\) is the stroke period and \(\tau\) is the drive time.  In the
simulation we take \(T = \SI{2.0}{s}\) and \(\tau = \SI{0.9}{s}\), giving a
drive‑to‑recovery ratio consistent with on‑water rowing.  The drive force
\(F\_\text{drive}\) is chosen as \SI{3000}{N}, corresponding to eight rowers
delivering roughly \SI{375}{N} each during the drive.

\subsection{Energy accounting}

The instantaneous mechanical power supplied by the rowers is \(P\_\text{in} =
F\_{\text{thrust}}\,v\).  The instantaneous power dissipated by drag is
\(P\_\text{drag} = F\_d\,v\).  Integrating these quantities over time yields
the total energy input and total energy lost to drag.  The change in boat
kinetic energy is \(\Delta E\_k = \tfrac{1}{2} m [v(t\_f)^2 - v(0)^2]\).
The mechanical efficiency of the stroke cycle is defined as

\begin{equation}
  \eta = \frac{\Delta E\_k}{E\_\text{input}},
  \label{eq:efficiency}
\end{equation}

representing the fraction of rower work that appears as an increase in
kinetic energy.

\section{Numerical Method}

Equation\,\eqref{eq:newton} is integrated numerically on a uniform time grid
with step \(\Delta t = \SI{0.01}{s}\).  At each step the thrust
\eqref{eq:thrust} is evaluated, the drag force is computed from the
previous velocity using \eqref{eq:drag}, and the acceleration is
calculated.  An explicit Euler update advances the velocity.  Non‑negativity
is enforced to avoid artificial reversal during recovery.  The power
quantities are accumulated at each time step to obtain energy input and
energy lost to drag.

The simulation runs for \SI{20}{s}, covering ten full strokes.  To
characterise steady behaviour, average velocity is computed over the second
half of the time span.  Figures are generated for boat velocity and thrust
as functions of time.

\section{Results}

\subsection{Speed and thrust profiles}

Figure\,\ref{fig:velocity} shows the boat speed over twenty seconds.
Starting from rest, the shell accelerates rapidly during the first few
strokes.  A periodic sawtooth pattern emerges: velocity increases during
each drive and decays during the recovery as the boat coasts under drag.
After about six seconds the system approaches a limit cycle with peak
speed near \SI{6.8}{m/s} and troughs near \SI{4.9}{m/s}.  The average
velocity over the final ten seconds is \num{5.785}\,\si{m/s},
corresponding to about \SI{20.8}{km/h}.

\begin{figure}[h]
  \centering
  \includegraphics[width=0.85\linewidth]{speed_plot.png}
  \caption{Boat velocity as a function of time.  The model predicts
    oscillatory motion, with rapid acceleration during the drive and
    deceleration during the recovery.  After several strokes the system
    settles into a limit cycle with an average speed around
    \SI{5.8}{m/s}.}
  \label{fig:velocity}
\end{figure}

Figure\,\ref{fig:thrust} depicts the thrust profile.  The square‑wave
character reflects the simplified stroke model: rowers apply a constant
force of \SI{3000}{N} for \SI{0.9}{s} followed by a recovery of
\SI{1.1}{s}.  This representation neglects finer details of blade
engagement and release but captures the alternating pattern of work and
coast.

\begin{figure}[h]
  \centering
  \includegraphics[width=0.85\linewidth]{thrust_plot.png}
  \caption{Thrust applied by the rowers over time.  The drive portion
    lasts \SI{0.9}{s} and the recovery \SI{1.1}{s}, producing a
    sawtooth speed profile as shown in Figure\,\ref{fig:velocity}.}
  \label{fig:thrust}
\end{figure}

\subsection{Energy and efficiency}

Table\,\ref{tab:summary} summarises the energy balance for the twenty‑second
simulation.  The rowers delivered roughly \num{1.41e5}\,\si{J} of mechanical
energy, of which only about \num{9.17e3}\,\si{J} became an increase in
kinetic energy.  The rest was dissipated by drag.  The resulting
mechanical efficiency is therefore only \SI{6.5}{\percent}.  The large
drag losses emphasise the importance of reducing hull resistance and
smoothing the stroke to minimise speed fluctuations.

\begin{table}[h]
  \centering
  \caption{Summary of simulation results (20 s, 10 strokes).}
  \label{tab:summary}
  \begin{tabular}{@{}lrr@{}}
    \toprule
    Quantity & Value & Unit \\
    \midrule
    Average speed (last 10 s) & 5.785 & m/s \\
    Maximum speed & 6.759 & m/s \\
    Total energy input & $1.41\times 10^5$ & J \\
    Drag energy loss & $1.32\times 10^5$ & J \\
    Increase in kinetic energy & $9.17\times 10^3$ & J \\
    Mechanical efficiency & 6.49 & \% \\
    \bottomrule
  \end{tabular}
\end{table}

\section{Discussion}

The simplified model captures key features of rowing dynamics: oscillatory
speed, rapid initial acceleration and a limit cycle speed reached after
several strokes.  The predicted average speed around \SI{5.8}{m/s}
aligns with typical racing paces for an eight‑person crew.  The efficiency
estimate of 6--7 \% is low compared with energy conversion efficiencies in
cycling or running, reflecting the significant hydrodynamic drag that must
be overcome.  Real rowing efficiency is highly dependent on stroke timing,
blade design and hull shape.  Our model uses a constant drag constant of
\SI{40}{N\,s^2/m^2} tuned to achieve realistic speeds; measured drag
coefficients can vary widely【779142687997843†L388-L450】 and may be lower
for well‑designed shells.  The square‑wave thrust model also ignores
variations in force during the drive and the dynamics of rower body mass
shifting within the boat.

Nevertheless, the framework illustrates how small changes in parameters
alter performance.  Increasing the drive force raises both peak and
average speeds but also accentuates speed oscillations.  Increasing the
drag coefficient dramatically reduces speed and improves efficiency only
marginally, because more thrust energy is expended fighting drag.  More
realistic models could include blade slip, lift forces, crew timing
differences and elastic energy stored in the boat.  Computational fluid
dynamics models explicitly calculate blade drag and lift, but their
complexity often obscures intuition.  The present approach provides a
transparent, adaptable platform for exploring energy transfer in rowing.

\section{Conclusion}

A simple one‑dimensional model combining periodic thrust and quadratic
hydrodynamic drag predicts realistic speeds for an eight‑person racing
shell and quantifies the flow of mechanical energy.  By integrating
Newton’s equation with an empirical drag law the model reproduces the
sawtooth speed pattern seen on the water and demonstrates that most of
the rowers’ work is dissipated by drag.  Despite its simplicity, the
framework highlights key trade‑offs in boat design and stroke strategy.
Future work could incorporate variable thrust profiles, rower body
dynamics and measured drag coefficients to refine the predictions.

\bibliographystyle{plain}
\begin{thebibliography}{9}

\bibitem{drag_equation}
  Drag equation.\ \emph{Wikipedia}.  The drag force on an object moving
  through a fluid is given by \(F\_d = \tfrac{1}{2}\rho u^2 c\_d A\), where
  \(\rho\) is fluid density, \(u\) velocity, \(A\) reference area and
  \(c\_d\) is the drag coefficient【443245353880283†L120-L134】.  Accessed
  6~Dec~2025.

\bibitem{racing_shell_drag}
  J.~G.~Buckmann and S.~D.~Harris.\ An experimental determination of the drag
  coefficient of a Mens~8+ racing shell.\ \emph{SpringerPlus}, 3:512, 2014.\
  The study measured masses for rowers (\SI{655}{kg}), shell (\SI{89}{kg}) and
  oars (\SI{22.5}{kg})【779142687997843†L270-L287】 and noted that the drag force
  on a racing shell can be expressed as a constant times speed to a power
  between 1.5 and 2【779142687997843†L256-L261】.  Average drag coefficients
  around 10 were reported across trials【779142687997843†L388-L450】.

\end{thebibliography}

\end{document}